% Conditional accounting lemma; ordinary tree/Kraft-mass counting, not a
% newly established structural bound for unknot recognition.
\subsection{A precise conditional global search budget}
\label{sec:conditional-global-budget}

\begin{theorem}[Search-tree budget]
\label{thm:search-tree-budget}
Let an algorithm on an input of size $n$ have a finite rooted search tree
$\mathcal T$, including rejected children, repeated states, and restarts.
Write $b(v)$ for the number of children of $v$, and let $H(n)$ bound the
number of edges on every root-to-leaf path.  Suppose that every such path
$P$ satisfies
\[
 \sum_{v\in P}\log\max\{1,b(v)\}\leq B(n).
\]
If all work is chargeable to tree nodes or incoming edges at cost at most
$W(n)$ per node, then
\[
 \#\operatorname{leaves}(\mathcal T)\leq e^{B(n)},\qquad
 |\mathcal T|\leq 1+H(n)e^{B(n)},
\]
and the total cost is at most
$W(n)\bigl(1+H(n)e^{B(n)}\bigr)$.
In particular, polynomial $H,W$ and $B(n)=O((\log n)^2)$ give
$n^{O(\log n)}$ time; $B(n)=O((\log n)^3)$ gives
$n^{O((\log n)^2)}$ time.
\end{theorem}

\begin{proof}
Give the root mass $1$, and divide the mass of every internal node equally
among its children.  Since the tree is finite, the leaf masses sum to $1$.
The mass of a leaf is the reciprocal of the product of the positive
branching numbers along its path, and is therefore at least $e^{-B(n)}$.
This proves the leaf bound.  Every nonroot node lies on a root-to-leaf
path, so the union of these paths contains at most $H(n)$ nonroot nodes
per leaf.  Summing the assumed node charges proves the cost bound.
\end{proof}

Here $W(n)$ includes state construction, exact arithmetic, surface tests,
certificate production and replay, and the generation of children.
Generating a large child list may instead be charged incrementally to
the respective child nodes.  Failed surface candidates and unsuccessful
LP attempts must either occur explicitly in the tree or be covered by
these local charges.  A sequence of unary restarts has zero branching
charge but consumes depth; the hypothesis on $H$ prevents it from being
free.  In particular, polynomial encoded state size alone does not bound
the work of the present simplex backends.

\begin{corollary}[Charging dimension decreases and their replenishment]
Suppose a nonnegative integer $D(v)$ is assigned to every tree node.
Designate certain binary branching nodes as conflict nodes, and assume
$D$ decreases by at least one on either outgoing edge of every such node.
For a path $P$, let
\[
 R(P)=\sum_{\substack{(v,w)\in P\\v\text{ not a conflict node}}}
       \max\{0,D(w)-D(v)\},\qquad
 C(P)=\sum_{\substack{v\in P\\v\text{ not a conflict node}}}
       \log\max\{1,b(v)\}.
\]
Then its branching budget is at most
\[
 (\log 2)\bigl(D(\mathrm{root})+R(P)\bigr)+C(P).
\]
\end{corollary}

\begin{proof}
If the path traverses $r$ conflict branches, telescoping the changes in
$D$ gives $0\leq D(\mathrm{leaf})\leq D(\mathrm{root})-r+R(P)$.
Each conflict node contributes $\log 2$; the other contributions sum to
$C(P)$.
\end{proof}

The intrinsic-dimension conflict theorem supplies the required decrease
within a fixed matching problem.  Replacing a sector, changing the source,
or restarting a hierarchy can replenish the dimension; the positive
increases must be charged through $R(P)$, or through a separate proved
phase budget.  This is a sufficient accounting target, not a bound on
$R(P)$, $C(P)$, or hierarchy depth.  To obtain a recognizer one must also
prove that the search rules cover every required possibility and that all
terminal conclusions are valid; neither the counting lemma nor a capped
or inconclusive leaf establishes those topological facts.
